\conclusions

Con el desarrollo del videojuego educativo \textit{Go Goals! Digital} se comprueba y reafirma el éxito y cumplimiento del propósito de esta investigación, obteniendo una solución a todas las problemáticas planteadas y abriendo la puerta a la innovación en el ámbito de los videojuegos educativos enfocados en la promoción de los Objetivos de Desarrollo Sostenible. El proceso de investigación y el análisis completo de la solución propuesta han permitido llegar a las siguientes conclusiones:

\begin{itemize}
    \item El marco teórico sobre videojuegos educativos, aprendizaje basado en juegos y gamificación demuestra el potencial pedagógico de estas herramientas en la educación para el desarrollo sostenible.
    
    \item El análisis de soluciones homólogas proporcionó patrones de diseño y mecánicas de juego que sirvieron como referencia en el desarrollo del sistema, identificando prácticas que debían evitarse para garantizar el éxito de la investigación.
    
    \item La definición de requisitos funcionales y no funcionales ayudó a establecer la estructura técnica y visual de la solución y garantizar una mayor comprensión del videojuego a desarrollar.
    
    \item El proceso de implementación da como resultado el sistema deseado que cumple con todas las especificaciones técnicas planteadas, además de contar con una estructura sólida desde su base metodológica hasta su arquitectura de desarrollo.
    
    \item Las pruebas realizadas después del desarrollo confirman que \textit{Go Goals! Digital} cumple con todas las expectativas y puede ser calificado como un videojuego educativo versátil, útil y necesario.
\end{itemize}